% !PRINTED 295
\weberpage{295}
%<<CONT
and leads to the congruences:
\weq{5}{\begin{aligned}
1.&\quad \alpha'\alpha+\beta'\gamma \equiv \gamma'\\
2.&\quad \alpha'\beta+\beta'\delta \equiv \delta'\\
3.&\quad \gamma'\alpha+\delta'\gamma \equiv -\alpha'+\gamma'\xi\\
4.&\quad \gamma'\beta+\delta'\delta \equiv -\beta'+\delta'\xi
\end{aligned}\ \ (\mathrm{mod}\ p),}
and from 1. and 2. there follows further:
\weq{6}{\begin{aligned}
5.&\quad \alpha' \equiv \gamma'\delta-\delta'\gamma\\
6.&\quad \beta' \equiv -\gamma'\beta+\delta'\alpha
\end{aligned}\ \ (\mathrm{mod}\ p).}
If these values are substituted in (5) 3. and 4., it follows that:
\[
\begin{aligned}
\gamma'(\alpha+\delta-\xi) &\equiv 0\\
\delta'(\alpha+\delta-\xi) &\equiv 0
\end{aligned}\ \ (\mathrm{mod}\ p),
\]
from which, since $\gamma'$, $\delta'$ cannot both be divisible by $p$, it follows that:
\weq{7}{\xi \equiv \alpha+\delta \quad (\mathrm{mod}\ p);}
if $\xi$ is determined in this way, then in (5) the congruences 3., 4. follow
from 1., 2. But if one substitutes $\gamma'$, $\delta'$ from 1., 2. in
\weq{8}{\alpha'\delta'-\beta'\gamma' \equiv 1 \quad (\mathrm{mod}\ p)}
then one obtains
\weq{9}{\alpha'^2\beta+\alpha'\beta'(\delta-\alpha)-\beta'^2\gamma \equiv 1 \quad (\mathrm{mod}\ p).}
If $\alpha'$, $\beta'$ are determined from this, then all the requirements stated
in (5), and hence also in (2), (4), are satisfied.

It therefore remains only to show that the congruence (9) is always
solvable.

This possibility is evident if $\beta$ or $-\gamma$ is a quadratic
residue of $p$; for then
\[
\begin{aligned}
\alpha' &\equiv \sqrt{\beta^{-1}}, & \beta' &\equiv 0\\
\text{or}\quad \alpha' &\equiv 0, & \beta' &\equiv \sqrt{-\gamma^{-1}}
\end{aligned}\ \ (\mathrm{mod}\ p)
\]
satisfies the requirement laid down.

In what follows we must now distinguish whether the substitution $S$
belongs to class I, II, III of the preceding section.

If first
\[
\text{I.} \qquad \left(\frac{\alpha+\delta}{2}\right)^2 \equiv 1 \quad (\mathrm{mod}\ p),
\]
then, if we put $S^m$ in place of $S$, by \S\,79, (9),
$\beta$ passes into $m\beta$, $\gamma$ into $m\gamma$, and $m$ can always be so determined
that $m\beta$ or $-m\gamma$ becomes a quadratic residue of $p$.

Here belongs also the case that $\beta \equiv 0$ or $\gamma \equiv 0$ and
$\alpha \equiv \delta \ (\mathrm{mod}\ p)$.
